\lecture{8}{Mon 19 Sep 2022 10:46}{}

\subsection{The homomorphism theorems}

\begin{theorem}
	If $N \triangleleft G$, then there exists a surjective homomorphism $\varphi: G \to  G \ N$ with $N = \operatorname{ker} \varphi$.
\end{theorem}

\begin{theorem}[First Homomorphism Theorem]
	(2.7.1)
	Let $\varphi: G \to  G'$ be a surjective homomorphism of groups with kernel $K$. Then $G' \cong G / K$.
\end{theorem}
\begin{proof}
	Define the map  \begin{align*}
		\psi: G / K &\longrightarrow G' \\
		Ka &\longmapsto \psi(Ka) = \varphi(k)
	.\end{align*}
	Check this gives an  isomorphism:

	\begin{itemize}
	\item (Well-defined property)
		We check that if $Ka = Kb$ then $\psi(Ka) = \psi(Kb)$. But $\psi(Ka) = \varphi(a)$ and  $\psi(Kb) = \varphi(b)$. Then if $Ka = Kb$, then $K = K b a^{-1} $, whence $ba^{-1}  \in K$, so that \[
			\varphi(b)\varphi(a)^{-1} = \varphi(ba^{-1} ) = e' \implies \varphi(b) = \varphi(a)
		.\] 
		Thus $\psi(Ka) = \psi(Kb)$.
	\item (Homomorphism)
		Whenever $a, b \in G$, \[
			\psi(Ka) \psi(Kb) 
			=\varphi(a)\varphi(b)
			=\varphi(ab)
			=\psi(Kab)
			=\psi(Ka \cdot Kb)
		.\] 
	\item ($\psi$ surjective)
		Whenever $g \in G'$, by the surjective property of $\varphi$, there exists $a \in G$ such that $\varphi(a) = g$, and $\psi(Ka) = \varphi(a)$, so there exists $Ka \in G / K$ such that $\psi(Ka) = g$.
	\item ($\psi$ injective)
		Whenever $\psi(Ka) = \psi(Kb)$, one has $\varphi(a) = \varphi(b)$, so $\varphi(ab^{-1} )=\varphi(a) \varphi(b)^{-1}= e'$, whence $ab^{-1} \in \operatorname{ker} \varphi = K$.
		Thus $a \in K b$, whence $Ka = Kb$. Thus $\psi$ injective.
	\end{itemize}
	So, $\psi$ is a well-defined bijective homomorphism from $G / K$ to $G'$, and thus $G / K \cong G'$.
\end{proof}

\begin{theorem}[Correspondence Theorem]
	(2.7.2)
	Let $\varphi: G \to  G'$ be a surjective homomorphism of groups with kernel $K$. Suppose $H' \le G'$ and \[
		H = \{a \in G: \varphi(a) \in H'\} 
	.\] 
	Then $H \le G$, $K \subset H$, and $H / K \cong H'$. Also, whenever $H' \triangleleft G'$ then $H \triangleleft G$.
\end{theorem}
\begin{remark}
	When $H' = \{e\} $, $H$ is the kernel. We have $\varphi$ condensing information from $H$ to $H'$, forgetting about the kernel. We can still work back and understand what is going on.
\end{remark}
\begin{proof}
	Check: 
	\begin{enumerate}
		\item $H\le G$ :
		We have $e \in H$, so $H \neq  \emptyset $. Also, whenever $a, b \in H$, one has $\varphi(ab^{-1} ) = \varphi(a)\varphi(b)^{-1} \in H'$, and hence
		$ab^{-1} \in H$. So $H \le G$.
	\end{enumerate}
	Then we use the first homomorphism theorem.
\end{proof}

